%%%%%%%%%%%%%%%%%%%%%%%%%%%%%%%%%%%%%%%%%%%%%%%%%
%%%%              CONCLUSIONS                %%%%
%%%%%%%%%%%%%%%%%%%%%%%%%%%%%%%%%%%%%%%%%%%%%%%%%

\section{Conclusions}
\label{sec:conclusions}

% ── 4a  ACHIEVEMENTS ──────────────────────────────────────────────────────────
\subsection{Phase 1 Achievements}
\label{sec:achievements}

We delivered a complete, modular IR pipeline on schedule (deadline Apr~13),
covering corpus loading, indexing, query building, retrieval, evaluation,
and result serialisation — all Colab-compatible.%
\footnote{Phase~1 notebook:
\url{https://colab.research.google.com/drive/1UvDhFXYEzZxse6gknDyBSUaq0EAHZkFV?usp=sharing}.}
Five retrieval strategies were implemented and compared: BM25,
LM-Jelinek-Mercer, LM-Dirichlet, dense KNN, and RRF fusion, following a
disciplined tuning protocol — query field ablation first, then 5-fold
cross-validation per strategy, with the test set touched exactly once.

The most impactful single design choice was selecting the MedCPT encoder,
which yielded $\Delta$NDCG@100 = +0.0501 on the test set
($0.7576 \to 0.8077$ versus the msmarco baseline; +0.0825 on 5-fold CV).
The locked Phase~2 strategy, KNN\,(MedCPT), reaches
NDCG@100 = 0.8077, MAP = 0.6143, and R@100 = 0.8933 on 33 test
topics; RRF (BM25+MedCPT, $k = 60$) pushes NDCG@100 to 0.8245 on the
same set.
The high Recall@100 of 0.8933 ensures a strong evidence pool for Phase~2
generation.

Several architectural decisions contributed to these results.
Our multi-field index design materialises every parameter configuration as a
dedicated OpenSearch field, enabling instant replay and side-by-side
comparison without re-encoding.
A systematic pooling-mode analysis (CLS vs.\ mean vs.\ mean-no-special)
established that architecture-appropriate pooling is critical before
comparing encoder families.
The exhaustive RRF sweep over 9~strategy pairs and 3~$k$ values confirmed
that sparse+dense fusion (BM25+KNN) is consistently the best pairing,
exploiting complementary lexical and semantic signals.
Finally, dual qrel scales (binary and graded) allowed evaluation with both
precision-oriented (MAP) and gain-discounted (NDCG) metrics, while the
robust odd/even topic split ensured no information leakage between selection
and final evaluation.

% ── 4b  LIMITATIONS ───────────────────────────────────────────────────────────
\subsection{Limitations}
\label{sec:limitations}

Several limitations should be noted.  The relevance judgements are derived
from automated BioGen submissions rather than human assessors, so systematic
biases are possible.  Moreover, the mean of 46.1 relevant documents per topic
is $3\text{--}15\times$ above standard TREC collections, which inflates MAP
and recall figures and makes them not directly comparable to published
benchmarks.  The 33-topic test set is relatively small, and the
fixed odd/even split offers no cross-validation guarantee that topic
difficulty is perfectly balanced.  A handful of relevant PMIDs referenced in
the qrels are absent from the 4{,}194-abstract corpus, creating an
irreducible recall ceiling.

Additionally, while the encoder comparison used brute-force cosine to control
for approximation noise, production KNN retrieval still relies on HNSW
indexing, which introduces some recall loss.

\subsection*{Future Work}

Phase~2 will add cross-encoder reranking with the
MedCPT-Cross-Encoder, LLM-augmented generation via vLLM with automatic
citation grounding, and LLM-as-judge evaluation.
Phase~3 will introduce a ReAct-style agentic loop — with Planner, Explorer,
Aggregator, and Report Writer roles — for multi-step reasoning over the
retrieval backbone.
Beyond the project scope, promising directions include learned sparse
retrieval (SPLADE) and late interaction (ColBERT) for higher recall at lower
rank cutoffs, as well as obtaining human relevance assessments to validate
the automated qrel quality.
